
\documentclass{beamer}

\usepackage[spanish]{babel}
\usepackage[utf8]{inputenc}
\usepackage[T1]{fontenc}
\usepackage{amsmath, amssymb}
\usepackage{graphicx}
\usepackage{mwe} % para 'example-image' como imagen de muestra

\title{semanal 2}
\author{KevimBr}
\date{22/08/2025}

\begin{document}

% --- DIAPOSITIVA 1 ---
\begin{frame}[fragile]
  \frametitle{Kevin Rafael Buenrostro Rodríguez}

\begin{figure}
  \centering
  \includegraphics[width=0.6\textwidth]{20241223_210050.jpg}
  \caption{mis perritos}
  \label{fig:mi-imagen}
\end{figure}


\end{frame}

% Segunda diapositiva (Teorema de Dirac)
\begin{frame}{Teorema de Dirac}
  \textbf{Enunciado:}  
  El Teorema de Dirac, en teoría de grafos, establece que un grafo simple y conexo con al menos tres vértices, 
  donde cada vértice tiene un grado mayor o igual a la mitad del número total de vértices, es hamiltoniano.  

  \medskip
  En otras palabras: si cada vértice tiene suficientes conexiones, entonces existe un ciclo que pasa por todos los vértices exactamente una vez.  

  \medskip
  \textit{Escogí este teorema porque me recordó al reto del "trazo con un solo trazo" que vimos en secundaria, 
  pero ahora con razones matemáticas.}

\cite{dirac1952}
\end{frame}


% Tercera diapositiva (resolver ecuación )

\begin{frame}{Resolviendo la ecuación}
  Resolver la ecuación:
  \[
     \frac{3}{4}x + 9 = 15
  \]

  \begin{align*}
     \frac{3}{4}x + 9 &= 15 \\[6pt]
     \frac{3}{4}x &= 15 - 9 \\[6pt]
     \frac{3}{4}x &= 6 \\[6pt]
     x &= \frac{6 \cdot 4}{3} \\[6pt]
     x &= 8
  \end{align*}

  \cite{stewart2015}
\end{frame}

% --- REFERENCIAS ---
\begin{frame}
  \frametitle{Referencias}
  \begin{thebibliography}{9}

  \bibitem{Dirac theorem}
    \textit Dirac, G.A. (1952). Some theorems on abstract graphs. 
    \emph{Proceedings of the London Mathematical Society}, 2(1), 69–81.

  \end{thebibliography}
\end{frame}

\bibliographystyle{plain}
\bibliography{referencias}


\end{document}